% Options for packages loaded elsewhere
% Options for packages loaded elsewhere
\PassOptionsToPackage{unicode}{hyperref}
\PassOptionsToPackage{hyphens}{url}
%
\documentclass[
  11pt,
  letterpaper,
  DIV=11,
  numbers=noendperiod,
  twoside]{scrartcl}
\usepackage{xcolor}
\usepackage[left=1.8in, right=1.8in, top=1.5in, bottom=1.5in,
includemp=TRUE, marginparwidth=0in, marginparsep=0in]{geometry}
\usepackage{amsmath,amssymb}
\setcounter{secnumdepth}{3}
\usepackage{iftex}
\ifPDFTeX
  \usepackage[T1]{fontenc}
  \usepackage[utf8]{inputenc}
  \usepackage{textcomp} % provide euro and other symbols
\else % if luatex or xetex
  \usepackage{unicode-math} % this also loads fontspec
  \defaultfontfeatures{Scale=MatchLowercase}
  \defaultfontfeatures[\rmfamily]{Ligatures=TeX,Scale=1}
\fi
\usepackage{lmodern}
\ifPDFTeX\else
  % xetex/luatex font selection
  \setmathfont[]{Garamond-Math}
\fi
% Use upquote if available, for straight quotes in verbatim environments
\IfFileExists{upquote.sty}{\usepackage{upquote}}{}
\IfFileExists{microtype.sty}{% use microtype if available
  \usepackage[]{microtype}
  \UseMicrotypeSet[protrusion]{basicmath} % disable protrusion for tt fonts
}{}
\usepackage{setspace}
% Make \paragraph and \subparagraph free-standing
\makeatletter
\ifx\paragraph\undefined\else
  \let\oldparagraph\paragraph
  \renewcommand{\paragraph}{
    \@ifstar
      \xxxParagraphStar
      \xxxParagraphNoStar
  }
  \newcommand{\xxxParagraphStar}[1]{\oldparagraph*{#1}\mbox{}}
  \newcommand{\xxxParagraphNoStar}[1]{\oldparagraph{#1}\mbox{}}
\fi
\ifx\subparagraph\undefined\else
  \let\oldsubparagraph\subparagraph
  \renewcommand{\subparagraph}{
    \@ifstar
      \xxxSubParagraphStar
      \xxxSubParagraphNoStar
  }
  \newcommand{\xxxSubParagraphStar}[1]{\oldsubparagraph*{#1}\mbox{}}
  \newcommand{\xxxSubParagraphNoStar}[1]{\oldsubparagraph{#1}\mbox{}}
\fi
\makeatother


\usepackage{longtable,booktabs,array}
\usepackage{calc} % for calculating minipage widths
% Correct order of tables after \paragraph or \subparagraph
\usepackage{etoolbox}
\makeatletter
\patchcmd\longtable{\par}{\if@noskipsec\mbox{}\fi\par}{}{}
\makeatother
% Allow footnotes in longtable head/foot
\IfFileExists{footnotehyper.sty}{\usepackage{footnotehyper}}{\usepackage{footnote}}
\makesavenoteenv{longtable}
\usepackage{graphicx}
\makeatletter
\newsavebox\pandoc@box
\newcommand*\pandocbounded[1]{% scales image to fit in text height/width
  \sbox\pandoc@box{#1}%
  \Gscale@div\@tempa{\textheight}{\dimexpr\ht\pandoc@box+\dp\pandoc@box\relax}%
  \Gscale@div\@tempb{\linewidth}{\wd\pandoc@box}%
  \ifdim\@tempb\p@<\@tempa\p@\let\@tempa\@tempb\fi% select the smaller of both
  \ifdim\@tempa\p@<\p@\scalebox{\@tempa}{\usebox\pandoc@box}%
  \else\usebox{\pandoc@box}%
  \fi%
}
% Set default figure placement to htbp
\def\fps@figure{htbp}
\makeatother


% definitions for citeproc citations
\NewDocumentCommand\citeproctext{}{}
\NewDocumentCommand\citeproc{mm}{%
  \begingroup\def\citeproctext{#2}\cite{#1}\endgroup}
\makeatletter
 % allow citations to break across lines
 \let\@cite@ofmt\@firstofone
 % avoid brackets around text for \cite:
 \def\@biblabel#1{}
 \def\@cite#1#2{{#1\if@tempswa , #2\fi}}
\makeatother
\newlength{\cslhangindent}
\setlength{\cslhangindent}{1.5em}
\newlength{\csllabelwidth}
\setlength{\csllabelwidth}{3em}
\newenvironment{CSLReferences}[2] % #1 hanging-indent, #2 entry-spacing
 {\begin{list}{}{%
  \setlength{\itemindent}{0pt}
  \setlength{\leftmargin}{0pt}
  \setlength{\parsep}{0pt}
  % turn on hanging indent if param 1 is 1
  \ifodd #1
   \setlength{\leftmargin}{\cslhangindent}
   \setlength{\itemindent}{-1\cslhangindent}
  \fi
  % set entry spacing
  \setlength{\itemsep}{#2\baselineskip}}}
 {\end{list}}
\usepackage{calc}
\newcommand{\CSLBlock}[1]{\hfill\break\parbox[t]{\linewidth}{\strut\ignorespaces#1\strut}}
\newcommand{\CSLLeftMargin}[1]{\parbox[t]{\csllabelwidth}{\strut#1\strut}}
\newcommand{\CSLRightInline}[1]{\parbox[t]{\linewidth - \csllabelwidth}{\strut#1\strut}}
\newcommand{\CSLIndent}[1]{\hspace{\cslhangindent}#1}



\setlength{\emergencystretch}{3em} % prevent overfull lines

\providecommand{\tightlist}{%
  \setlength{\itemsep}{0pt}\setlength{\parskip}{0pt}}



 


\setlength\heavyrulewidth{0ex}
\setlength\lightrulewidth{0ex}
\usepackage[automark]{scrlayer-scrpage}
\clearpairofpagestyles
\cehead{
  Brian Weatherson
  }
\cohead{
  New Work for Dempster-Shafer Functions
  }
\ohead{\bfseries \pagemark}
\cfoot{}
\makeatletter
\newcommand*\NoIndentAfterEnv[1]{%
  \AfterEndEnvironment{#1}{\par\@afterindentfalse\@afterheading}}
\makeatother
\NoIndentAfterEnv{itemize}
\NoIndentAfterEnv{enumerate}
\NoIndentAfterEnv{description}
\NoIndentAfterEnv{quote}
\NoIndentAfterEnv{equation}
\NoIndentAfterEnv{longtable}
\NoIndentAfterEnv{abstract}
\renewenvironment{abstract}
 {\vspace{-1.25cm}
 \quotation\small\noindent\emph{Abstract}:}
 {\endquotation}
\newfontfamily\tfont{EB Garamond}
\addtokomafont{disposition}{\rmfamily}
\addtokomafont{title}{\normalfont\itshape}
\let\footnoterule\relax

\makeatletter
\renewcommand{\@maketitle}{%
  \newpage
  \null
  \vskip 2em%
  \begin{center}%
  \let \footnote \thanks
    {\itshape\huge\@title \par}%
    \vskip 0.5em%  % Reduced from default
    {\large
      \lineskip 0.3em%  % Reduced from default 0.5em
      \begin{tabular}[t]{c}%
        \@author
      \end{tabular}\par}%
    \vskip 0.5em%  % Reduced from default
    {\large \@date}%
  \end{center}%
  \par
  }
\makeatother
\RequirePackage{lettrine}

\renewenvironment{abstract}
 {\quotation\small\noindent\emph{Abstract}:}
 {\endquotation\vspace{-0.02cm}}

\setmainfont{EB Garamond Math}[
  BoldFont = {EB Garamond SemiBold},
  ItalicFont = {EB Garamond Italic},
  RawFeature = {+smcp},
]

\newfontfamily\scfont{EB Garamond Regular}[RawFeature=+smcp]
\renewcommand{\textsc}[1]{{\scfont #1}}

\renewcommand{\LettrineTextFont}{\scfont}
\newcommand{\nmodels}{\mathrel{\ooalign{$\models$\cr\raisebox{-0.001ex}{\hss$\mkern-1mu/\hss$}}}}
\newcommand{\llbracket}{[\![}
\newcommand{\rrbracket}{]\!]}
\setlength\heavyrulewidth{0ex}
\setlength\lightrulewidth{0ex}
\usepackage[lines=2]{lettrine}
\cehead{Draft of February 28, 2026}
\DeclareSymbolFont{symbolsC}{U}{txsyc}{m}{n}
\DeclareMathSymbol{\boxright}{\mathrel}{symbolsC}{128}
\cehead{
       Brian Weatherson
       }
\KOMAoption{captions}{tableheading}
\makeatletter
\@ifpackageloaded{caption}{}{\usepackage{caption}}
\AtBeginDocument{%
\ifdefined\contentsname
  \renewcommand*\contentsname{Table of contents}
\else
  \newcommand\contentsname{Table of contents}
\fi
\ifdefined\listfigurename
  \renewcommand*\listfigurename{List of Figures}
\else
  \newcommand\listfigurename{List of Figures}
\fi
\ifdefined\listtablename
  \renewcommand*\listtablename{List of Tables}
\else
  \newcommand\listtablename{List of Tables}
\fi
\ifdefined\figurename
  \renewcommand*\figurename{Figure}
\else
  \newcommand\figurename{Figure}
\fi
\ifdefined\tablename
  \renewcommand*\tablename{Table}
\else
  \newcommand\tablename{Table}
\fi
}
\@ifpackageloaded{float}{}{\usepackage{float}}
\floatstyle{ruled}
\@ifundefined{c@chapter}{\newfloat{codelisting}{h}{lop}}{\newfloat{codelisting}{h}{lop}[chapter]}
\floatname{codelisting}{Listing}
\newcommand*\listoflistings{\listof{codelisting}{List of Listings}}
\makeatother
\makeatletter
\makeatother
\makeatletter
\@ifpackageloaded{caption}{}{\usepackage{caption}}
\@ifpackageloaded{subcaption}{}{\usepackage{subcaption}}
\makeatother
\usepackage{bookmark}
\IfFileExists{xurl.sty}{\usepackage{xurl}}{} % add URL line breaks if available
\urlstyle{same}
\hypersetup{
  pdftitle={New Work for Dempster-Shafer Functions},
  pdfauthor={Brian Weatherson},
  hidelinks,
  pdfcreator={LaTeX via pandoc}}


\title{New Work for Dempster-Shafer Functions}
\author{Brian Weatherson}
\date{2026}
\begin{document}
\maketitle
\begin{abstract}
TBD
\end{abstract}


\setstretch{1.1}
Conventional wisdom these day is that credences should be probabilistic,
and this is because of accuracy considerations. I'm going to argue that
this conventional wisdom relies on an assumption that we have reason to
reject. The assumption is that we should be \emph{realists} about
anything we have credences about, in the sense that there is a fact of
the matter about anything we should have a credence about. If we drop
this assumption, if we even drop a priori certainty in there being a
fact of the matter, it could be coherent for credences to not be
probabilistic. The right coherence constraint for someone who is unsure
about realism is not probability, but Dempster-Shafer belief functions.
This isn't what belief functions were designed for; there isn't much
evidence that the people who designed them intended to question realism.
But it is work, new work, that these belief functions can do.

To motivate the idea that we should be uncertain about realism, I'll
work through some questions about fiction. The particular questions
aren't that important in themselves, in fact part of the point is that
they aren't important, but they are questions where anti-realism seems
plausible.

For example, how many cousins does Sherlock Holmes have?\footnote{The
  usual example, following Lewis (\citeproc{ref-Lewis1978b}{1978}),
  involves how many hairs Holmes has. I prefer this question because it
  makes clear the issue is not vagueness.} Does he have an odd or even
number of cousins? If you're like me, untutored intuition says that (1)
is true, while there is no fact of the matter about (2) and (3).

\begin{enumerate}
\def\labelenumi{(\arabic{enumi})}
\tightlist
\item
  Holmes has either an odd or even number of cousins.
\item
  Holmes has an odd number of cousins.
\item
  Holmes has an even number of cousins.
\end{enumerate}

Intuition isn't the most reliable guide around here, and there are well
known complications with trying to validate all these intuitions.
Nevertheless, I'll argue these intuitions are correct. At least, I'll
argue that the latter intuition is correct; for this paper I'll simply
assume that (1) is true, and argue that there is no fact of the matter
about (2) and (3).

In doing so, I'll be arguing against the orthodox view represented by
David Lewis (\citeproc{ref-Lewis1978b}{1978}). His view is that
sentences like (2) and (3) have an implicit operator, roughly \emph{In
the Holmes fiction}, and this operator is a normal necessity
operator.\footnote{Lewis's account of that operator uses counterfactual
  conditionals, but on the account of counterfactuals he endorses
  (\citeproc{ref-Lewis1973a}{Lewis 1973}), \emph{If it were that A} is a
  normal necessity operator, so we get back what I said in the text.} So
just like \(\Box (A \vee B)\) can be true while \(\Box A\) and
\(\Box B\) are both false, Lewis says that (1) can be true while (2) and
(3) are both false.

Recently, Kyle Blumberg and Ben Holguín
(\citeproc{ref-BlumbergHolguin2025}{2025}) have argued that Lewis's
account has implausible predictions about fictional claims embedded in
attitude ascriptions. For instance, it implies that (4) is true.

\begin{enumerate}
\def\labelenumi{(\arabic{enumi})}
\setcounter{enumi}{3}
\tightlist
\item
  I know whether Holmes has either an odd or even number of cousins.
\end{enumerate}

They argue that Lewis's view should be rejected in favour of their
\emph{fictional reality} view. On their view, the Holmes stories pick
out a particular possible world, though it is unknowable which one. So
one of (2) and (3) is true. This seems incredible; it's more plausible
that there is no fact of the matter about how many cousins Holmes has.

What does it even mean to say there is `no fact of the matter', and is
this even coherent? I'll answer these questions by appeal to
\emph{possibility semantics}, as developed by Lloyd I. L. Humberstone
(\citeproc{ref-Humberstone1981}{1981}). In possibility semantics, the
role of possible worlds is played by \emph{possibilities}, which can be
incomplete. Despite this incompleteness, the logic of truth preservation
across possibilities is more or less classical.

I'm going to argue that a story does not pick out a possible world, but
instead some Humberstone-style possibilities. Which ones does a
particular story pick out? It's coherent to be unsure about that. Even
with the text in front of them, a reader might not know the
extra-textual information that settles some facts in the story. In such
a situation, their credences about what's true in a story might be
distributed across various posisbilities.

At this point we get (back) to belief functions. Measures on possibility
frames do not always determine probability functions; rather, they might
determine Dempster-Shafer belief functions. So, I'll conclude, it is
rational for someone uncertain about what is true in a story to have
credences that track one of these belief functions.

This conclusion might seem surprising. There is a famous argument,
originally due to James Joyce (\citeproc{ref-Joyce1998}{1998}), that any
non-probabilistic credence function is accuracy dominated. This argument
seems to rule out exactly the view I put forward in the last paragraph.
That argument, however, assumes exactly the kind of realism that is
implausible in the case of thoughts about fiction. Someone who rejects
fictional realism can happily use Dempster-Shafer belief functions
without worrying about being excessively inaccurate.

Indeed, this last point generalises to any case where realism isn't
clearly true. I'll suggest this opens up a response on behalf of the
dogmatist to an important objection to dogmatism due to Roger White
(\citeproc{ref-White2006}{2006}). (What I say here builds on what James
Pryor (\citeproc{ref-Pryor2007}{2007}) says in an unpublished but widely
discussed note.)

So this paper lives at the intersection of possibility semantics,
accuracy dominance arguments for probabilism, and theories of truth in
fiction. The paper will work its way forward along this list of topics
then back again. So Section~\ref{sec-humberstone} and
Section~\ref{sec-classical} discuss possibility semantics. I'll start by
setting out the details of them possibility semantics, and end by
discussing the question, which will have become somewhat pressing by
then, of whether the theory is really \emph{classical}.
Section~\ref{sec-joyce} and Section~\ref{sec-realism} will discuss
accuracy dominance arguments, first discussing why accuracy dominance
arguments are so important, and then coming back to why the anti-realist
can reject them. And in the middle, in Section~\ref{sec-fiction}, I'll
endorse Blumberg and Holguìn's objections to Lewis's theory of truth in
fiction, while rejecting their alternative in favour of an anti-realist
model.

\section{Possibility Semantics}\label{sec-humberstone}

A Humberstone possibility frame (hereafter, just \emph{possibility
frame}) is a triple \(\langle W, \leqslant, \mathbb{R}, f, V\rangle\),
where \(W\) is a set of possibilities, \(\leqslant\) is a
\emph{refinement} relation on them, \(\mathbb{R}\) is a countable family
of accesibility relations (with \(R_i\) as the \(i\)'the member), \(f\)
is a selection function, and \(V\) is a valuation function. I'll set
aside \(\mathbb{R}\) and \(f\) until later, since they only play a role
in interpreting modal statements and counterfactual conditionals
respectively. I'll focus first on \(W, \leqslant\) and \(V\).

One way to think about the possibilities in \(W\), suggested by I. L.
Humberstone (\citeproc{ref-Humberstone1981}{1981}) and used extensively
by Wesley Holliday (\citeproc{ref-Holliday2025}{2025}), is that they are
sets of possible worlds. I prefer to think of them as \emph{stories}, at
least on the assumption that stories are incomplete.\footnote{This
  assumption is part of what's at issue in the paper, but I'm not
  begging any questions here. Whether one pictures possibilities as
  stories or sets of worlds does not bear any argumentative weight.} A
possibility, like a story, is something which settles some facts, while
possibly leaving others open.

Following Humberstone, I'll write \(x \leqslant y\) to mean that \(y\)
is a \textbf{refinement} of \(x\).\footnote{\(y \geqslant x\) will just
  mean the same thing as \(x \leqslant y\).} To say \(y\) refines \(x\)
is to say that \(y\) makes true everything that \(x\) makes true, and
possibly some more things besides. If possibilities are sets of possible
worlds, then \(y\) is a subset of \(x\). (This makes the notation here
confusing, which is why Holliday (\citeproc{ref-Holliday2025}{2025})
prefers to write \(y \sqsubseteq x\) when \(y\) is a refinement of
\(x\).) If possibilities are stories, then \(y\) is a refinement of
\(x\) if \(y\) fills in part of what \(x\) leaves open. This is how
I'll, intuitively, think of refinements.

Finally, \(V\) is a \textbf{valuation} function. It tells us where the
atomics are true. Formally, it is a function from \(\mathcal{P}\), the
set of atomics, to \(\wp(W)\), , satisfying these two constraints:

\begin{itemize}
\tightlist
\item
  For all \(x\), if \(x \in V(p)\) and \(x \leqslant y\), then
  \(y \in V(p)\). Intuitively, truth for atomics is \textbf{persistent}
  across refinements.
\item
  For all \(x\), if
  \(\forall y \geqslant x\;\, \exists z \geqslant y: z \in V(p)\), then
  \(x \in V(p)\). This is what L. Humberstone
  (\citeproc{ref-Humberstone2011}{2011, 900}) calls
  \textbf{refinability}. It means that \(p\) only fails to be true at
  \(x\) if there is some refinement of \(x\) where it is settled as
  being untrue.
\end{itemize}

Humberstone then proposes the the following theory of truth at a
possibility for the language of propositional logic. (We'll treat
\(\rightarrow\) as a defined connective, with
\(A \rightarrow B =_{df} \neg A \vee B\).)

\begin{align*}
[\text{Vbls}] \quad & \mathcal{M} \models_x p_i \text{ iff } x \in V(p_i); \\
[\neg] \quad & \mathcal{M} \models_x \neg A \text{ iff } \forall y \geqslant x, \, \mathcal{M} \nmodels_y  A; \\
[\wedge] \quad & \mathcal{M} \models_x A \wedge B \text{ iff } \mathcal{M} \models_x A \text{ and } \mathcal{M} \models_x B; \\
[\vee] \quad & \mathcal{M} \models_x A \vee B \text{ iff } \forall y \geqslant x \;\, \exists z \geqslant y : \mathcal{M} \models_z A \text{ or } \mathcal{M} \models_z B. \\
[\rightarrow] \quad & \mathcal{M} \models_x A \rightarrow B \text{ iff } \mathcal{M} \models_x \neg(A \wedge \neg B)
\end{align*}

The first two clauses are familiar, and the third looks much like the
intuitionist treatment of negation. The final clause amounts to treating
\(\rightarrow\) as a defined connective. The surprising one is
\([\vee]\). It says that a disjunction is true iff at any refinement, at
least one of the disjuncts could become true. This is not the most
obvious understanding of \(\vee\), but it allows Humberstone to prove
three striking results.

First, every sentence in the language is persistent. If
\(\mathcal{M} \models_x A\) and \(x \leqslant y\), then
\(\mathcal{M} \models_y A\). For any sentence, truth is always preserved
when moving to a refinement.

Second, refinability holds for all sentences in the language. Once we've
added negation to the language, we can restate refinability as the
condition that, for any \(A\), if \(\mathcal{M} \nmodels_x A\), there is
some refinement \(y\) of \(x\) such that
\(\mathcal{M} \models_y \neg A\). Intuitively, if a sentence is not
true, then at some refinement its negation is true.

The third point is the most important. Say that \(\Gamma \vDash_{pf} A\)
iff at any point in any possibility frame where all the members of
\(\Gamma\) are true, \(A\) is also true. And say that
\(\Gamma \vdash_{CL} A\) if \(\Gamma\) classically entails \(A\). Then
Humberstone shows that all of the canonical inference rules for
\(\vdash_{CL}\) also hold for \(\vDash_{pf}\), and hence \(\vdash_{CL}\)
and \(\vDash_{pf}\) exactly overlap.

It's useful to work through one of these rules to show how the system
works. Because it will play a role in what follows, I'll show that
disjunction elimination holds just as one would expect. One natural
formulation\footnote{There are other formulations which make weaker
  structural assumptions, but this one will do to show that we're
  preserving classical logic.} of this is that if
\(\Gamma \cup \{A\} \vdash_{CL} C\),
\(\Delta \cup \{B\} \vdash_{CL} C\), and \(\Phi \vdash_{CL} A \vee B\)
then \(\Gamma \cup \Delta \cup Phi \vdash_{CL} C\). Let's show that the
same principle holds for \(\vDash_{pf}\).

Assume that \(\Gamma \cup \{A\} \vDash_{pf} C\),
\(\Delta \cup \{B\} \vDash_{pf} C\), and \(\Phi \vDash_{pf} A \vee B\).
We want to show that \(\Gamma \cup \Delta \cup \Phi \vDash_{pf} C\). So
take any possibility \(x\) where all members of
\(\Gamma \cup \Delta \cup \Phi\) are true, and suppose for reductio that
\(\mathcal{M} \nmodels_x C\). By refinability, there is some
\(y \geqslant x\) such that \(\mathcal{M} \models_y \neg C\), i.e., for
all \(z \geqslant y\), \(\mathcal{M} \nmodels_z C\). By persistence, all
members of \(\Gamma \cup \Delta \cup \Phi\) are true at \(y\). Since all
of \(\Phi\) is true at \(y\), we have
\(\mathcal{M} \models_y A \vee B\). By \([\vee]\), there is some
\(z \geqslant y\) such that \(\mathcal{M} \models_z A\) or
\(\mathcal{M} \models_z B\). By persistence again, all members of
\(\Gamma \cup \Delta\) are true at \(z\). If
\(\mathcal{M} \models_z A\), then all of \(\Gamma \cup \{A\}\) are true
at \(z\), so \(\mathcal{M} \models_z C\). If
\(\mathcal{M} \models_z B\), then all of \(\Delta \cup \{B\}\) is true
at \(z\), so \(\mathcal{M} \models_z C\). Either way,
\(\mathcal{M} \models_z C\). But \(z \geqslant y\) and
\(\mathcal{M} \models_y \neg C\), so \(\mathcal{M} \nmodels_z C\).
Contradiction.

Humberstone's system looks non-classical. It includes points that make
true neither \(A\) nor \(\neg A\) for some \(A\). It includes a rule for
negation that looks more like intuitionistic negation than classical
negation. But the system is still classical. All the sequents, and all
the admissible rules on sequents, carry over from classical logic. We
have a non-standard semantics for a very standard logic.

Possibility frames have attracted some attention recently as an
interesting classical alternative to possible worlds semantics.
Humberstone's original application of them was to modal logic, and
Wesley Holliday (\citeproc{ref-Holliday2025}{2025}) has written a
thorough study of that application, differing in some key respects from
Humberstone's original picture.\footnote{The differences are largely
  about the relationship between the \(R\)-relation in modal semantics
  and \(\leqslant\), so they aren't relevant to our story.} Matthew
Harrison-Trainor (\citeproc{ref-HarrisonTrainor2019}{2019}) has shown
how to extent the semantics to first-order logic. Brian Weatherson
(\citeproc{ref-WeathersonHoPF}{2026}) applies possibility frame
semantics to a puzzle about counterfactuals. And Fabrizio Cariani
(\citeproc{ref-CarianiUnpub}{2025}) uses possibility frames to model the
open future. This paper extends that list; I'm going to use them to
provide a better way of talking about fictional truth.

Since many theories of fiction involve modals or counterfactuals, we
need to say more about \(\mathbb{R}\) and \(f\). First we'll talk about
modals.

In order to make sure persistence and refinability hold in the modal
part of the language, Humberstone puts some constraints on what can go
into \(\mathbb{R}\). (As noted above, Holliday prefers weaker
constraints, but the difference won't be important here.) Any
accessibility relation \(R_i\) has to satisfy these
constraints.\footnote{The names are taken from Holliday; Humberstone
  uses less descriptive names.}

\begin{description}
\tightlist
\item[\textbf{UpR}:]
If \(x \leqslant x'\) and \(x' R_i y\), then \(x R_i y\).
\item[\textbf{RDown}:]
If \(x R_i y\) and \(y \leqslant y'\), then \(x R_i y'\).
\item[\textbf{RRef++}:]
If \(x R_i y\), then there exists \(x' \geqslant x\) such that for all
\(x'' \geqslant x'\), \(x'' R_i y\).
\end{description}

\textbf{UpR} says that if a refinement of \(x\) can access \(y\), then
\(x\) itself can already access \(y\): accessibility is not something
that can be gained by adding detail to the source. \textbf{RDown} is a
converse of this; it says that accessibility cannot be lost by adding
detail to the target. \textbf{RRef++} says that if \(x\) can access
\(y\), there is some refinement \(x'\) of \(x\) where it is settled that
\(x'\) can access \emph{y}. This last access can't be overturned by
further refinement of \(x'\).

Given these constraints, the truth condition for the box operator is:

\begin{align*}
[\Box_i] \quad & \mathcal{M} \models_x \Box_i A \text{ iff } \forall y \, (x R_i y \Rightarrow \mathcal{M} \models_y A)
\end{align*}

This should be familiar: \(\Box_i A\) is true at \(x\) iff \(A\) is true
at every \(R_i\)-accessible possibility.

Humberstone treats \(\Diamond\) as a defined connective,
\(\Diamond_i A\) just means \(\neg \Box_i \neg A\), and I'll do the
same. It is easy to show that if \(R_i\) satisfies these three
constraints, then even after adding \(\Box_i\) and \(\Diamond_i\) to the
language, all sentences will satisfy persistence and refinability.

Weatherson (\citeproc{ref-WeathersonHoPF}{2026}) describes a theory of
counterfactuals in possibility semantics that validates conditional
excluded middle. That's what we're going to need, but unfortunately that
semantics doesn't guarantee persistence and refinability. Here is a
modified version of his semantics that does.

We posit a selection function \(f\) as in the semantics for
counterfactuals suggested by Stalnaker
(\citeproc{ref-Stalnaker1968}{1968}). We assume that the selection
function follows the these constraints.

\begin{description}
\tightlist
\item[\textbf{Conditional Excluded Middle}]
Either \(f(A, x) = \emptyset\), or
\(\exists y: f(A, x) = \{z : z \geqslant y\}\).
\item[\textbf{Truth}]
\(f(A, x) \subseteq A\)
\item[\textbf{Centering}]
\(x \in f(A, x)\) whenever \(x \in A\)
\item[\textbf{Nearness}]
If \(f(A, x) \subseteq B\) and \(f(B, x) \subseteq A\), then
\(f(A, x) = f(B, x)\)
\item[\textbf{UpR for \(f\)}]
If \(x \leqslant x'\) and \(y \in f(A, x')\), then \(y \in f(A, x)\)
\item[\textbf{RDown for \(f\)}]
If \(y \in f(A, x)\) and \(y \leqslant y'\), then \(y' \in f(A, x)\)
\item[\textbf{RRef++ for \(f\)}]
If \(y \in f(A, x)\), then there exists \(x' \geqslant x\) such that for
all \(x'' \geqslant x'\), \(y \in f(A, x'')\)
\end{description}

Then the truth condition for counterfactuals is:

\begin{align*}
[\boxright] \quad & \mathcal{M} \models_x A \boxright B \text{ iff } \forall y \in f(A, x) \, \mathcal{M} \models_y B 
\end{align*}

I'll say that \(f(A, x)\) are the nearest \(A\)-worlds; it will become
clear why this is a reasonable locution.

\textbf{Conditional Excluded Middle} says that there are no nearest
\(A\)-worlds, or they are all and only the refinements of some
particular possibility. So it can't be that neither \(B\) nor \(\neg B\)
holds throughout the nearest \(A\)-possibilities: if there is a nearest
\(A\)-possibility \(y\), then by refinability for \(B\), either
\(\mathcal{M} \models_y B\) (in which case persistence gives
\(\mathcal{M} \models_z B\) for all \(z \geqslant y\), so
\(A \boxright B\) is true), or there is some \(y' \geqslant y\) with
\(\mathcal{M} \models_{y'} \neg B\) (in which case persistence gives
\(\mathcal{M} \models_z \neg B\) for all \(z \geqslant y'\), and since
\(\{z : z \geqslant y'\} \subseteq \{z : z \geqslant y\} = f(A, x)\),
there is a nearest \(A\)-possibility where \(\neg B\) holds). This
validates Conditional Excluded Middle for counterfactuals:
\(\mathcal{M} \models_x (A \boxright B) \vee (A \boxright \neg B)\).

The other principles are more straightforward. \textbf{Truth} says that
the nearest \(A\)-worlds have to actually make \(A\) true.
\textbf{Centering} says that if \(A\) is true at \(x\), then \(x\) is
one of the nearest \(A\)-worlds. This system doesn't quite have Strong
Centering, since there could be other nearby \(A\)-worlds, but it is
close. The nearest \(A\)-worlds to an \(A\)-world can only be
coarsenings or refinements of it. \textbf{Nearness} says that if the
nearest \(A\)-worlds are all \(B\)-worlds, and the nearest \(B\)-worlds
are all \(A\)-worlds, then the nearest \(A\)-worlds just are the nearest
\(B\)-worlds. That has to be true if \(f\) is really measuring nearness.
Finally, the last three clauses ensure that if we define \(R_A\) such
that \(xR_Ay\) iff \(y \in f(A, x)\), then \(R_A\) is an acceptable
accessibility relation. This in turn guarantees that the language is
still persistent and refinable, even with \(\boxright\) added.

\section{Accuracy Dominance}\label{sec-joyce}

Over the twentieth century, there were several important pragmatic
arguments for probabilism, the thesis that credences are coherent only
if they are probabilistic. The most important of these were Dutch Book
arguments and representation theorem arguments. These arguments were
problematic in two respects. First, they were \emph{pragmatic}
arguments. Second, they were arguments \emph{for probabilism}. The first
problem was that these seemed to show at most that non-probabilistic
credences were pragmatically sub-optimal, not that they violated a
distinctive doxastic norm. The second problem was that these arguments
didn't engage, at all, with non-probabilistic alternatives. It's worth
pausing on this because it brings up something rather distinctive about
the literature on probabilism.

Most philosophical arguments engage at length, some would say inordinate
length, with rival views. Pick up a work arguing for this or that
functionalist thesis in philosophy of mind, and you expect to see
extensive treatment of the flaws of non-functionalist rivals. And that's
the normal case. Sometimes, a philosopher will put forward a purportedly
sound deductive for a distinctive conclusion, as Frank Jackson
(\citeproc{ref-Jackson1982}{1982}) did with the knowledge argument.
Those arguments don't turns on comparative considerations weighing up,
in this case, dualism vs physicalism. But in those cases, the arguments
will be carefully crafted to be maximally compelling to people tempted
by the opposing view.

The arguments for probabilism have none of that. Some premises are
suggested, and they are shown to entail probabilism. There is no
consideration of which of these premises might be more compelling to
someone tempted by this or that non-probabilist view. Indeed, in most of
these papers there is no sense of which non-probabilistic view the
author is most concerned to argue against. Or, in many cases, whether
there is \emph{any} non-probabilistic view that the author thinks is
worth careful consideration.

The risk of this approach is that the arguments will be
question-begging. As Patrick Maher (\citeproc{ref-Maher1997}{1997})
showed, that's exactly what happened. While there are many ways to not
be a probabilist, the most widely discussed in the modern literature is
that credences should be Dempster-Shafer belief functions (hereafter,
simply \emph{belief functions}). Maher showed that these arguments for
probabilism were question-begging against the view that credences should
be belief functions.

David Christensen (\citeproc{ref-Christensen1996}{1996}) provided the
first important attempt to deal with the first problem, i.e., that these
arguments showed that non-probabilistic credences violated the wrong
kind of norm. He argued that the non-probabilist was committed to
valuing the same thing two different ways under different modes of
presentation, and this was a genuinely doxastic incoherence. This does,
I think, deal with the first problem, but the cost is that it makes the
second problem worse. A key premise of Christensen's argument is that
the value of a book of bets which can't both win is the sum of the
values of each bet. And there is no reason for the belief functions
theorist to accept that. The argument is just clearly question-begging
against them.\footnote{This premise is also rejected by some opponents
  of expected value maximisation, like Lara Buchak
  (\citeproc{ref-BuchakRisk}{2013}).}

This is the

\section{Fictional Anti-Realism}\label{sec-fiction}

\begin{itemize}
\tightlist
\item
  Agreement with Lewis and with BH - hidden operator
\item
  Agreement with BH against Lewis - the wonder cases are problems
\item
  Disagreement with BH - realism is literally incredible
\item
  First, appeal to intuition - especially about songs
\item
  Second, appeal to possibility of consistent but incompatible fanfics
\item
  Conclusion, what's true in a story is what's true in a
  \emph{possibility}
\end{itemize}

\section{Accuracy for Anti-Realists}\label{sec-realism}

\begin{itemize}
\tightlist
\item
  Start with 1/3 p, 1/3 \textasciitilde p
\item
  Accuracy dominance says in any world, you'll be better with 1/2, 1/2
\item
  Illustrate with Brier score
\item
  But in a possibility, you're better with 1/3, 1/3
\item
  In general, any belief function is optimal wrt some possibility frame
\item
  Conclusion: The accuracy dominance argument assumes realism, and
  anti-realists can reject it. Belief functions are credence functions
  for anti-realists
\item
  Compare Horowitz's guessing framework. Even clearer that the realist
  assumption is being made.
\item
  Objection 1: This doesn't save Arthur Dempster and Glenn Shafer. Sure,
  it doesn't.
\item
  Objection 2: This rests on a strange way of understanding disjunction.
  Sure, but disjunction is strange. We shouldn't have a hundred times as
  many books and papers on conditionals as on disjunction. This is why
  possibilities are in the disjunction chapter of \emph{The
  Connectives}.
\item
  Objection 3: This is only relevant to fiction, and to trivial
  questions about fiction.
\item
  First, it's not just trivial. Any hard question about a fiction should
  have a possibility that there is no fact of the matter, so some weight
  to a non-world possibility, so non-probabilistic credences.
\item
  Second, it's not just fiction. It's anything where realism might fail.
\item
  In low evidence situations, that might include \emph{anything}. We
  might not be external world realists with no evidence.
\item
  That gives us a response to White on dogmatism
\item
  Objection 4. This doesn't explain the unreasonable effectiveness of
  probabilism
\item
  Reply. Probabilism given realism. In many cases, realism is a safe
  assumption. Compare Priest on how classical logic is the right logic
  for situations known to be consistent. I say the same thing, except
  for me it's situations known to be complete.
\end{itemize}

\section{Is This Still Classical Logic?}\label{sec-classical}

\begin{itemize}
\tightlist
\item
  What is a logic?
\item
  Option 1. A set of theorems. PF preserves the theorems
\item
  Option 2. A set of axioms plus rules for generating theorems. Still
  good, we get the axioms and we get modus ponens.
\item
  Option 3. A set of single-conclusion sequents. We get that, as shown
  in Section~\ref{sec-humberstone}
\item
  Option 4. A set of basic sequents (i.e., \(\Gamma \vdash A\) whenever
  \(A \in \Gamma\)) plus rules for generating more. We get that too, as
  also shown. This distinguishes possibility semantics from
  supervaluationism. The supervaluationist does not get, e.g.,
  conditional proof, as Williamson showed.
\item
  Option 5. A set of multiple-conclusion sequents. We do not get that.
  Hence we don't get option 6; basic multiple conclusion sequents plus
  rules
\item
  Include a pair of 3 by 2 tables, showing what possibility semantics
  preserves from classical logic, and what supervaluationism preserves,
  with one key difference
\item
  But what to say about the last row, with multiple conclusions
\item
  On the one hand, multiple conclusion logic is weird; cite Dummett,
  Wright, etc on how weird it is
\item
  On the other hand, that's really what folks have commonly assumed.
  It's a common part of accuracy based arguments that one or other
  complete state is true. Inferring that from zero premises is a
  multiple-conclusion argument.
\item
  So there is a sense in which this is a non-classical response to
  accuracy dominance. But I think it's a kind of non-classicality which
  we have reason to accept at least some of the time.
\end{itemize}

\section*{References}\label{references}
\addcontentsline{toc}{section}{References}

\phantomsection\label{refs}
\begin{CSLReferences}{1}{0}
\bibitem[\citeproctext]{ref-BlumbergHolguin2025}
Blumberg, Kyle, and Ben Holguín. 2025. {``Fictional Reality.''}
\emph{Philosophical Review} 134 (2): 149--201. doi:
\href{https://doi.org/10.1215/00318108-11676103}{10.1215/00318108-11676103}.

\bibitem[\citeproctext]{ref-BuchakRisk}
Buchak, Lara. 2013. \emph{Risk and Rationality}. Oxford: Oxford
University Press.

\bibitem[\citeproctext]{ref-CarianiUnpub}
Cariani, Fabrizio. 2025. {``Modeling Future Indeterminacy in
Possibilitysemantics.''} \url{https://philarchive.org/rec/CARMFI}.

\bibitem[\citeproctext]{ref-Christensen1996}
Christensen, David. 1996. {``Dutch-Book Arguments {D}e-Pragmatized:
Epistemic Consistency for Partial Believers.''} \emph{Journal of
Philosophy} 93 (9): 450--79. doi:
\href{https://doi.org/10.2307/2940893}{10.2307/2940893}.

\bibitem[\citeproctext]{ref-HarrisonTrainor2019}
Harrison-Trainor, Matthew. 2019. {``First-Order Possibility Models and
Finitary Completeness Proofs.''} \emph{Review of Symbolic Logic} 12 (4):
637--62. doi:
\href{https://doi.org/10.1017/S1755020319000418}{10.1017/S1755020319000418}.

\bibitem[\citeproctext]{ref-Holliday2025}
Holliday, Wesley H. 2025. {``Possibility Frames and Forcing for Modal
Logic.''} \emph{Australasian Journal of Logic} 22 (2): 44--288. doi:
\href{https://doi.org/10.26686/ajl.v22i2.5680}{10.26686/ajl.v22i2.5680}.

\bibitem[\citeproctext]{ref-Humberstone1981}
Humberstone, I. L. 1981. {``From Worlds to Possibilities.''}
\emph{Journal of Philosophical Logic} 10 (3): 313--39. doi:
\href{https://doi.org/10.1007/BF00293423}{10.1007/BF00293423}.

\bibitem[\citeproctext]{ref-Humberstone2011}
Humberstone, Lloyd. 2011. \emph{The Connectives}. Cambridge, MA: MIT
Press.

\bibitem[\citeproctext]{ref-Jackson1982}
Jackson, Frank. 1982. {``Epiphenomenal Qualia.''} \emph{The
Philosophical Quarterly} 32 (127): 127--36. doi:
\href{https://doi.org/10.2307/2960077}{10.2307/2960077}.

\bibitem[\citeproctext]{ref-Joyce1998}
Joyce, James M. 1998. {``A Non-Pragmatic Vindication of Probabilism.''}
\emph{Philosophy of Science} 65 (4): 575--603. doi:
\href{https://doi.org/10.1086/392661}{10.1086/392661}.

\bibitem[\citeproctext]{ref-Lewis1973a}
Lewis, David. 1973. \emph{Counterfactuals}. Oxford: Blackwell
Publishers.

\bibitem[\citeproctext]{ref-Lewis1978b}
---------. 1978. {``Truth in Fiction.''} \emph{American Philosophical
Quarterly} 15 (1): 37--46. Reprinted in his \emph{Philosophical Papers},
Volume 1, Oxford: Oxford University Press, 1983, 261-275. References to
reprint.

\bibitem[\citeproctext]{ref-Maher1997}
Maher, Patrick. 1997. {``Depragmatised Dutch Book Arguments.''}
\emph{Philosophy of Science} 64 (2): 291--305. doi:
\href{https://doi.org/10.1086/392552}{10.1086/392552}.

\bibitem[\citeproctext]{ref-Pryor2007}
Pryor, James. 2007. {``Uncertainty and Undermining.''}
\url{https://www.jimpryor.net/research/papers/Uncertainty.pdf}.

\bibitem[\citeproctext]{ref-Stalnaker1968}
Stalnaker, Robert. 1968. {``A Theory of Conditionals.''} In
\emph{Studies in Logical Theory}, edited by Nicholas Rescher, 98--112.
Oxford: Blackwell.

\bibitem[\citeproctext]{ref-WeathersonHoPF}
Weatherson, Brian. 2026. {``Humberstone on Possibility Frames.''}
\url{https://brian.weatherson.org/quarto/posts/possibility-frames/HoPF.html}.

\bibitem[\citeproctext]{ref-White2006}
White, Roger. 2006. {``Problems for Dogmatism.''} \emph{Philosophical
Studies} 131 (3): 525--57. doi:
\href{https://doi.org/10.1007/s11098-004-7487-9}{10.1007/s11098-004-7487-9}.

\end{CSLReferences}



\noindent Draft for discussion.


\end{document}
